\section{The Debt of Nature}

The Catechism says what death is: ``In that `departure' which is death the soul is separated from the body.''\endnote{\emph{Catechism of the Catholic Church} 1005, official English (Vatican web text).} It rests on the definitions of two councils. The Council of Vienne defined that the rational soul is of itself and essentially the form of the human body, and the Fifth Lateran Council that this soul is immortal and is multiplied with the bodies into which it is infused.\endnote{Council of Vienne, constitution \latin{Fidei catholicae} (1312), DS 902; Fifth Lateran Council, session VIII, bull \latin{Apostolici regiminis} (1513), DS 1440.} A man is therefore one living thing, soul and body together, and not a spirit lodged in a body to be let out at the end like a bird from a cage. Death undoes that unity. The man dies; his spiritual soul survives, still knowing and loving, and waits in its separation for the body God will give back to it on the last day. That the souls of the just, once purified, see God even before they take up their bodies again is the teaching of Benedict~XII, expounded in \emph{What Death Does Not Finish}.\endnote{Benedict XII, constitution \latin{Benedictus Deus} (1336), DS 1000: \latin{etiam ante resumptionem suorum corporum}.}

Because body and soul belong together, their parting is violent, and the Fathers say so. Augustine will not let the faithful call it gentle: ``As regards bodily death, that is, the separation of the soul from the body, it is good unto none while it is being endured by those whom we say are in the article of death. For the very violence with which body and soul are wrenched asunder, which in the living had been conjoined and closely intertwined, brings with it a harsh experience, jarring horridly on nature so long as it continues\ldots''\endnote{Augustine, \emph{De civitate Dei} XIII.6, trans. Marcus Dods (Edinburgh: T. \& T. Clark, 1871).} Sometimes, he adds, the anguish is forestalled by a sudden stroke. However it comes, ``when it is piously and faithfully borne, it increases the merit of patience, but does not make the name of punishment inapplicable.''\endnote{\emph{De civitate Dei} XIII.6, Dods.} Whatever grace may do with death, then, it is in itself a wrench, a tearing apart of what God joined, and to feel it so is no failure of faith.

\subsection{Natural, and the wages of sin}

The Catechism gives two reasons why there should be such a tearing at all: ``In a sense bodily death is natural, but for faith it is in fact `the wages of sin.'\,''\endnote{CCC 1006, quoting Romans 6:23.} It is natural because man is made of the dust, a body composed of elements that, left to themselves, come apart. It is the wages of sin because God did not leave man to himself. Scripture says so in two places: ``For God made not death, neither hath he pleasure in the destruction of the living,'' and ``the wages of sin is death.''\endnote{Wisdom 1:13; Romans 6:23; Douay--Rheims (Challoner).} The Council of Trent defines that Adam, when he transgressed the commandment in paradise, lost the holiness and justice in which he had been set and incurred the death with which God had threatened him.\endnote{Council of Trent, session V (1546), decree on original sin, canon 1, DS 1511.}

Thomas Aquinas shows how the two answers fit together with an image from Aristotle. In his account sin causes death not directly, as a blow causes a wound, but as the man who pulls away a prop causes what was resting on it to fall: ``thus it is stated in \emph{Phys.} viii, text. 32, that `by displacing a pillar a man moves accidentally the stone resting thereon.' In this way the sin of our first parent is the cause of death and all such like defects in human nature, in so far as by the sin of our first parent original justice was taken away.''\endnote{Thomas Aquinas, \emph{Summa theologiae} I-II, q.~85, a.~5, corpus, trans. Fathers of the English Dominican Province.} The stone was always heavy; God's gift of original justice held it up, and when the gift was lost the stone fell, as it falls still on every child of Adam. The doctrine therefore concerns a gift: man was made for a life that the dust of which he is formed could not of itself secure, and he lost that life by turning from the Giver.

Ambrose, returning on the seventh day to the grave of his brother Satyrus, reads the same sentence of Genesis as a mercy. He points out that God did not say to Adam, ``because thou hast hearkened to the voice of the woman thou shalt return to the earth.'' The curse fell on the ground and its thorns. The return to the earth was named as the term of man's labour: ``In sweat shalt thou eat thy bread until thou return to the earth.'' ``Death is given for a remedy,'' Ambrose concludes, ``because it is the end of evils. \ldots{} You see that death is rather the goal of our penalties, by which an end is put to the course of this life.''\endnote{Ambrose, \emph{De excessu fratris sui Satyri} II.38, trans. H. de Romestin, in \emph{Nicene and Post-Nicene Fathers}, second series, vol.~10.} Augustine stresses the penalty and Ambrose the remedy, and the Church keeps both.

\subsection{The debt paid}

The tradition also calls death a debt. The anonymous \emph{Craft of Dying} of the fifteenth century counts among the blessings of a good death the ``paying of [the] debt of natural duty.''\endnote{\emph{The Book of the Craft of Dying}, ch.~1, ed. Frances M. M. Comper (London: Longmans, Green, 1917), brackets as printed.} The Church's prayer at the deathbed uses the same word as the priest commends the dying man to God: ``I commend thee to Almighty God, dearly beloved brother, and commit thee to Him whose creature thou art; that, when thou shalt have paid the debt of humanity by passing through death, thou mayest return to thy Maker, who formed thee from the dust of the earth.''\endnote{\emph{Rituale Romanum}, Ordo commendationis animae, the prayer \latin{Commendo te}, in \emph{A Manual of Prayers for the Use of the Catholic Laity} (New York: Catholic Publication Society, 1889), 516; Latin \latin{humanitatis debitum}.} The word fits both of the Church's answers: the debt is owed by nature, since every man is dust, and by justice, since every man is a son of Adam, and it is paid only once.

Scripture calls death an enemy, and the last to be destroyed.\endnote{1 Corinthians 15:26, Douay--Rheims (Challoner): ``And the enemy, death, shall be destroyed last.''} Until it is destroyed it keeps the violence Augustine describes. But it is not the end of everything, as the pagan takes it to be, because the debt that Adam contracted has been met by another, and death itself has been made to serve.
